\section{Finite-dimensional linear representations}

Let \(k\) be a field and let \(G\) be a finite group.

\begin{definition}[Finite-dimensional \(k\)-linear representation]
    A finite dimensional \(k\)-linear representation of \(G\) is a pair
    \[
        (M,\rho),
    \]
    where \(M\) is a finite-dimensional \(k\)-vector space and
    \[
        \rho:G\longrightarrow \operatorname{GL}(M)
    \]
    is a group homomorphism. Here
    \[
        \operatorname{GL}(M)
        =
        \{\,f:M\to M\mid f \text{ is an invertible } k\text{-linear map}\,\}.
    \]
\end{definition}

For \(g\in G\) and \(m\in M\), we often write
\[
    g\cdot m:=\rho(g)(m).
\]
Thus the representation condition says
\[
    (gh)\cdot m=g\cdot(h\cdot m),
    \qquad
    1_G\cdot m=m.
\]

\begin{definition}[The category of \(k\)-linear representations]
    The category of finite-dimensional \(k\)-linear representations of \(G\) is
    denoted by
    \[
        \operatorname{Rep}(G,\operatorname{vect}_k).
    \]
    Its objects are finite-dimensional \(k\)-linear representations
    \[
        (M,\rho).
    \]
    A morphism
    \[
        f:(M,\rho)\longrightarrow (N,\sigma)
    \]
    is a \(k\)-linear map
    \[
        f:M\longrightarrow N
    \]
    such that
    \[
        f(g\cdot m)=g\cdot f(m)
    \]
    for every \(m\in M\) and \(g\in G\). Equivalently,
    \[
        \sigma(g)\circ f=f\circ \rho(g),
        \qquad g\in G.
    \]
\end{definition}

\begin{definition}[Degree]
    Let \((M,\rho)\) be a finite-dimensional \(k\)-linear representation of \(G\).
    The dimension of \(M\) is called the degree of the representation:
    \[
        \deg(M,\rho):=\dim_k M.
    \]
\end{definition}

\begin{example}[The trivial representation]
    The trivial representation, also called the principal representation, is the
    one-dimensional representation
    \[
        \rho_{\operatorname{triv}}:G\longrightarrow \operatorname{GL}_1(k)
        \cong k^\times
    \]
    given by
    \[
        \rho_{\operatorname{triv}}(g)=1
    \]
    for all \(g\in G\). Equivalently, \(G\) acts on the one-dimensional vector
    space \(k\) by
    \[
        g\cdot x=x,
        \qquad g\in G,\ x\in k.
    \]
\end{example}

\begin{example}[One-dimensional representations]
    If \((M,\rho)\) is a \(k\)-linear representation with
    \[
        \dim_k M=1,
    \]
    then after choosing a basis of \(M\), we may identify
    \[
        \operatorname{GL}(M)\cong k^\times.
    \]
    Hence a one-dimensional representation is the same as a group homomorphism
    \[
        \rho:G\longrightarrow k^\times.
    \]
\end{example}

Let \((M,\rho)\) be a finite-dimensional \(k\)-linear representation of \(G\), and
assume
\[
    r:=\dim_k M.
\]
After fixing a basis of \(M\), we obtain an isomorphism of groups
\[
    \operatorname{GL}(M)\cong \operatorname{GL}_r(k),
\]
where
\[
    \operatorname{GL}_r(k)
\]
is the group of invertible \(r\times r\) matrices over \(k\).

Therefore a \(k\)-linear representation of \(G\) may be represented by a pair
\[
    (r,\rho),
\]
where
\[
    r\in \mathbb{Z}_{\geq 1}
\]
and
\[
    \rho:G\longrightarrow \operatorname{GL}_r(k)
\]
is a group homomorphism. Such a pair is called a matrix representation of \(G\)
over \(k\).

\begin{remark}
    Choosing a basis of \(M\) transforms each linear automorphism
    \[
        \rho(g):M\longrightarrow M
    \]
    into an invertible matrix. Thus a representation may be viewed concretely as
    assigning to every \(g\in G\) an invertible matrix \(\rho(g)\in
    \operatorname{GL}_r(k)\), satisfying
    \[
        \rho(gh)=\rho(g)\rho(h),
        \qquad
        \rho(1_G)=I_r.
    \]
\end{remark}

\begin{proposition}
    Let
    \[
        (r,\rho)
        \qquad\text{and}\qquad
        (s,\sigma)
    \]
    be two matrix representations of \(G\). Then they are isomorphic if and only
    if
    \[
        r=s
    \]
    and there exists a matrix
    \[
        A\in \operatorname{GL}_r(k)
    \]
    such that
    \[
        \sigma(g)=A\rho(g)A^{-1}
    \]
    for every \(g\in G\).
\end{proposition}

\begin{proof}
    Suppose first that the two representations are isomorphic. Then there exists
    a \(k\)-linear isomorphism
    \[
        f:k^r\longrightarrow k^s
    \]
    such that
    \[
        \sigma(g)\circ f=f\circ \rho(g)
    \]
    for every \(g\in G\). Since \(f\) is an isomorphism, we must have
    \[
        r=s.
    \]
    Let \(A\in \operatorname{GL}_r(k)\) be the matrix of \(f\) with respect to the
    chosen bases. Then the equivariance condition becomes
    \[
        \sigma(g)A=A\rho(g).
    \]
    Multiplying on the right by \(A^{-1}\), we obtain
    \[
        \sigma(g)=A\rho(g)A^{-1}.
    \]

    Conversely, suppose that \(r=s\) and that there exists
    \(A\in \operatorname{GL}_r(k)\) such that
    \[
        \sigma(g)=A\rho(g)A^{-1}
    \]
    for all \(g\in G\). Let
    \[
        f:k^r\longrightarrow k^r
    \]
    be the linear isomorphism whose matrix is \(A\). Then
    \[
        \sigma(g)A=A\rho(g),
    \]
    which means precisely that
    \[
        \sigma(g)\circ f=f\circ \rho(g)
    \]
    for every \(g\in G\). Hence \(f\) is an isomorphism of representations.
\end{proof}

\begin{remark}
    Thus two matrix representations \((r,\rho)\) and \((r,\sigma)\) are
    isomorphic precisely when the matrices
    \[
        \rho(g)
        \qquad\text{and}\qquad
        \sigma(g)
    \]
    are uniformly similar, meaning that the same change-of-basis matrix
    \(A\in \operatorname{GL}_r(k)\) satisfies
    \[
        \sigma(g)=A\rho(g)A^{-1}
    \]
    for every \(g\in G\).
\end{remark}

Let \(k\subseteq k'\) be a field extension, and let \(G\) be a group.

\begin{definition}[Extension of scalars]
    Let \((M,\rho)\) be a finite-dimensional \(k\)-linear representation of \(G\),
    where
    \[
        \rho:G\longrightarrow \operatorname{GL}(M).
    \]
    The extension of scalars of \((M,\rho)\) from \(k\) to \(k'\) is the
    \(k'\)-linear representation
    \[
        \bigl(k'\otimes_k M,\; k'\otimes_k \rho\bigr),
    \]
    where
    \[
        k'\otimes_k \rho:G\longrightarrow
        \operatorname{GL}_{k'}(k'\otimes_k M)
    \]
    is defined by
    \[
        g\longmapsto 1_{k'}\otimes \rho(g).
    \]
    Explicitly,
    \[
        (1_{k'}\otimes \rho(g))(\lambda\otimes m)
        =
        \lambda\otimes \rho(g)(m),
        \qquad
        \lambda\in k',\ m\in M.
    \]
\end{definition}

\begin{remark}
    If
    \[
        e_1,\dots,e_n
    \]
    is a \(k\)-basis of \(M\), then
    \[
        1\otimes e_1,\dots,1\otimes e_n
    \]
    is a \(k'\)-basis of \(k'\otimes_k M\). In particular,
    \[
        \dim_{k'}(k'\otimes_k M)=\dim_k M.
    \]
\end{remark}

\begin{remark}[Matrix form]
    Suppose that \((M,\rho)\) is given by a matrix representation
    \[
        \rho:G\longrightarrow \operatorname{GL}_n(k).
    \]
    Extending scalars to \(k'\) amounts to viewing each matrix
    \[
        \rho(g)\in \operatorname{GL}_n(k)
    \]
    as a matrix with entries in \(k'\):
    \[
        \rho(g)\in \operatorname{GL}_n(k')
    \]
    via the inclusion \(k\subseteq k'\).
\end{remark}

\begin{definition}[Representations rational over a subfield]
    Let \((M',\rho')\) be a \(k'\)-linear representation of \(G\). We say that
    \((M',\rho')\) is rational over \(k\) if there exists a \(k\)-linear
    representation \((M,\rho)\) of \(G\) such that
    \[
        (M',\rho')
        \cong
        \bigl(k'\otimes_k M,\; k'\otimes_k \rho\bigr)
    \]
    as \(k'\)-linear representations of \(G\).
\end{definition}

\begin{proposition}[Matrix criterion for rationality]
    Let
    \[
        \rho:G\longrightarrow \operatorname{GL}_n(k')
    \]
    be a matrix representation of \(G\) over \(k'\). Then \(\rho\) is rational
    over \(k\) if and only if there exists a matrix
    \[
        U\in \operatorname{GL}_n(k')
    \]
    such that, for every \(g\in G\), all entries of
    \[
        U\rho(g)U^{-1}
    \]
    belong to \(k\).
\end{proposition}

\begin{proof}
    The representation \(\rho\) is rational over \(k\) precisely when it is
    isomorphic over \(k'\) to the scalar extension of some representation
    \[
        \rho_0:G\longrightarrow \operatorname{GL}_n(k).
    \]
    An isomorphism of matrix representations over \(k'\) is given by a change of
    basis matrix \(U\in \operatorname{GL}_n(k')\), and the isomorphism condition is
    exactly
    \[
        U\rho(g)U^{-1}=\rho_0(g)
    \]
    for every \(g\in G\). Since \(\rho_0(g)\) has entries in \(k\), this proves one
    direction.

    Conversely, if such a matrix \(U\) exists, then the matrices
    \[
        U\rho(g)U^{-1}
    \]
    define a representation of \(G\) over \(k\), and \(\rho\) is obtained from it by
    extension of scalars to \(k'\). Hence \(\rho\) is rational over \(k\).
\end{proof}

Let
\[
    (M,\rho),\qquad (N,\sigma)
\]
be two \(k\)-linear representations of \(G\).

\begin{definition}[Direct sum]
    The direct sum of \((M,\rho)\) and \((N,\sigma)\) is the representation
    \[
        (M,\rho)\oplus (N,\sigma)
        :=
        (M\oplus N,\rho\oplus\sigma),
    \]
    where
    \[
        (\rho\oplus\sigma)(g)(m,n)
        =
        (\rho(g)m,\sigma(g)n).
    \]
\end{definition}

\begin{definition}[Tensor product]
    The tensor product of \((M,\rho)\) and \((N,\sigma)\) is the representation
    \[
        (M,\rho)\otimes (N,\sigma)
        :=
        (M\otimes_k N,\rho\otimes\sigma),
    \]
    where
    \[
        (\rho\otimes\sigma)(g)(m\otimes n)
        =
        \rho(g)m\otimes \sigma(g)n.
    \]
\end{definition}

\begin{definition}[Subrepresentation]
    Let \((M,\rho)\) be a \(k\)-linear representation of \(G\). A \(k\)-subspace
    \[
        M_1\subseteq M
    \]
    is called \(G\)-stable if
    \[
        \rho(g)(M_1)\subseteq M_1
    \]
    for every \(g\in G\).

    If \(M_1\subseteq M\) is \(G\)-stable, then the restricted maps
    \[
        \rho_1(g):=\rho(g)|_{M_1}
    \]
    define a representation
    \[
        (M_1,\rho_1),
    \]
    called a subrepresentation of \((M,\rho)\).
\end{definition}

\begin{remark}
    Since each \(\rho(g)\) is invertible and
    \[
        \rho(g)^{-1}=\rho(g^{-1}),
    \]
    the condition
    \[
        \rho(g)(M_1)\subseteq M_1
    \]
    for all \(g\in G\) implies
    \[
        \rho(g)(M_1)=M_1
    \]
    for all \(g\in G\).
\end{remark}

\begin{definition}[Quotient representation]
    Let \((M,\rho)\) be a \(k\)-linear representation of \(G\), and let
    \(M_1\subseteq M\) be a \(G\)-stable subspace. Then the quotient vector space
    \[
        M/M_1
    \]
    carries a natural representation
    \[
        \overline{\rho}:G\longrightarrow \operatorname{GL}(M/M_1)
    \]
    defined by
    \[
        \overline{\rho}(g)(m+M_1)
        :=
        \rho(g)(m)+M_1.
    \]
    This is called the quotient representation of \((M,\rho)\) by
    \((M_1,\rho_1)\).
\end{definition}

\begin{definition}[Simple representation]
    A \(k\)-linear representation \((M,\rho)\) of \(G\) is called simple, or
    irreducible, if its only \(G\)-stable subspaces are
    \[
        0
        \qquad\text{and}\qquad
        M.
    \]
\end{definition}

\begin{definition}[Indecomposable representation]
    A \(k\)-linear representation of \(G\) is called indecomposable if
    whenever
    \[
        (M,\rho)\cong
        (M_1,\rho_1)\oplus (M_2,\rho_2),
    \]
    one has
    \[
        M_1=0
        \qquad\text{or}\qquad
        M_2=0.
    \]
\end{definition}

\begin{remark}
    Every simple representation is indecomposable. The converse need not hold in
    general.
\end{remark}

\begin{theorem}[Krull--Schmidt theorem for finite-dimensional representations]
    Let \(G\) be a finite group, let \(k\) be a field, and let
    \[
        (M,\rho)
    \]
    be a finite-dimensional \(k\)-linear representation of \(G\).

    \begin{enumerate}
        \item The representation \((M,\rho)\) is a finite direct sum of
        indecomposable representations.

        \item Suppose that
        \[
            (M,\rho)
            \cong
            \bigoplus_{i\in I}(M_i,\rho_i)
            \cong
            \bigoplus_{j\in J}(M'_j,\rho'_j),
        \]
        where every \((M_i,\rho_i)\) and every \((M'_j,\rho'_j)\) is
        indecomposable. Then there exists a bijection
        \[
            \alpha:I\longrightarrow J
        \]
        such that, for every \(i\in I\),
        \[
            (M_i,\rho_i)\cong (M'_{\alpha(i)},\rho'_{\alpha(i)}).
        \]
    \end{enumerate}
\end{theorem}

\begin{proof}
    We regard \(k\)-linear representations of \(G\) as modules over the group
    algebra \(kG\). More precisely, a representation
    \[
        \rho:G\longrightarrow \operatorname{GL}(M)
    \]
    determines a left \(kG\)-module structure on \(M\) by
    \[
        \left(\sum_{g\in G} a_g g\right)m
        :=
        \sum_{g\in G} a_g\,\rho(g)(m),
        \qquad a_g\in k,\ m\in M.
    \]
    Conversely, every left \(kG\)-module gives a \(k\)-linear representation of
    \(G\). Hence \(\operatorname{Rep}(G,\operatorname{vect}_k)\) is equivalent to
    the category of finite-dimensional left \(kG\)-modules.

    Since \(G\) is finite, the group algebra \(kG\) is finite-dimensional over
    \(k\). Therefore every finite-dimensional \(kG\)-module has finite
    composition length.

    Applying the Krull--Schmidt theorem for finite-length modules to the
    \(kG\)-module \(M\), we obtain a decomposition
    \[
        M\cong M_1\oplus\cdots\oplus M_r
    \]
    into indecomposable \(kG\)-modules. Translating this back into
    representations gives a decomposition
    \[
        (M,\rho)
        \cong
        (M_1,\rho_1)\oplus\cdots\oplus(M_r,\rho_r)
    \]
    into indecomposable representations. This proves \textup{(1)}.

    For \textup{(2)}, suppose that
    \[
        (M,\rho)
        \cong
        \bigoplus_{i\in I}(M_i,\rho_i)
        \cong
        \bigoplus_{j\in J}(M'_j,\rho'_j),
    \]
    where all summands are indecomposable. Since \(M\) is finite-dimensional,
    both index sets \(I\) and \(J\) are finite.

    Passing again to \(kG\)-modules, this gives two decompositions of the same
    finite-length \(kG\)-module \(M\) into indecomposable direct summands:
    \[
        M\cong \bigoplus_{i\in I}M_i
        \cong
        \bigoplus_{j\in J}M'_j.
    \]
    By the Krull--Schmidt uniqueness theorem for finite-length modules, there
    exists a bijection
    \[
        \alpha:I\longrightarrow J
    \]
    such that
    \[
        M_i\cong M'_{\alpha(i)}
    \]
    as \(kG\)-modules for every \(i\in I\).

    Translating the \(kG\)-module isomorphisms back into 
    \(k\)-linear isomorphisms, we obtain
    \[
        (M_i,\rho_i)
        \cong
        (M'_{\alpha(i)},\rho'_{\alpha(i)})
    \]
    as \(k\)-linear representations of \(G\). This proves the desired uniqueness.
\end{proof}

\begin{remark}
    The theorem says that a finite-dimensional representation can be decomposed
    into indecomposable pieces, and that these pieces are unique up to
    isomorphism and reordering.
\end{remark}
